% !TEX root = ../notes.tex

\documentclass[../notes.tex]{subfiles}

\begin{document}

\section{May 11}
Let's summarize where we are.

\subsection{Idempotent Algebras}
In classical algebra, the fundamental object is the initial ring $\ZZ$. One frequently understands $\ZZ$ by its localizations, which look are $\QQ$ and $\ZZ_{(p)}$ for the various primes $p$. It is difficult to make these localizations directly make sense in higher algebra, but one can instead say that these localizations are the idempotent algebras $R$: namely, one has $R\otimes_\ZZ R=R$.
\begin{nex}
	The $\ZZ$-algebra $\FF_p$ has $\FF_p\otimes_\ZZ\FF_p=\FF_p\oplus\Sigma\FF_p$.
\end{nex}
While we're here, we note that one can already see some lattice structure among these idempotent algebras because $\ZZ_{(p)}$ naturally embeds into $\QQ$.

This course was about what is known for $\mathbb S$. Certainly one still has the idempotent algebras $\QQ$ and $\mathbb S_{(p)}$ for each prime $p$, but one could have more. They are found ``above'' $\mathbb S_{(p)}$.
\begin{theorem}[Thick subcategory] \label{thm:thick-subcat}
	Thick subcategories of perfect $\mathbb S_{(p)}$-modules are classified by type, where we say that a perfect subcategory of $\mathbb S_{(p)}$-modules has type at least $n$ if and only if the modules $K(m)\otimes X$ all vanish for $m\in\{0,1,\ldots,m-1\}$.
\end{theorem}
\begin{definition}
	We define the functor $L_n^f$ to be the localization of $\mathbb S_{(p)}$ obtained by formally killing colimits of type at least $n+1$ perfect $\mathbb S_{(p)}$-modules.
\end{definition}
\begin{remark}
	Note that $L_n^f\mathbb S_{(p)}$ is an idempotent algebra because it arises from some perfect modules.
\end{remark}
One could hope that all idempotent algebras should be found by such formal killing of certain thick subcategories of perfect modules, and this process only produces the idempotent $\mathbb S$-algebras $L_n^f\mathbb S$ for $n\ge1$. (This hope turns out to be false!)

Here are some other idempotent algebras.
\begin{definition}
	We define the idempotent algebra $L_n\mathbb S$ which is obtained by formally killing all spectra which are killed by $\{K(0),K(1),\ldots,K(n)\}$.
\end{definition}
\begin{remark}
	For $n\ge2$, it is known that $L_n\ne L_n^f$.
\end{remark}
\begin{remark}
	It turns out that all idempotent algebras lies between $L_n^f\mathbb S$ and $L_n\mathbb S$ for some $n$. The idea was to use \Cref{thm:thick-subcat} and the classification of skew fields.
\end{remark}
\begin{remark}
	Namely, there was a category of synthetic spectra
	\[\mathrm{SynSp}=\op{Mod}\lim\tau_{\ge2*}\op{BP}^{\otimes(\bullet+1)}.\]
	These are filtered objects, and the associated graded produces derived objects in $\op{QCoh}(\mc M_{\mathrm{fg},(p)})$. Taking the completion along the associated graded kills the difference between $L_n$ and $L_n^f$.
\end{remark}

\subsection{Hopkins's Conjecture}
One can hope to not care about the distinction between $L_n$ and $L_n^f$, which we will do for the rest of our time.
\begin{theorem}
	The limit of the diagram
	\[\mathbb S\to\mathbb S_{(p)}\to L_2\mathbb S\to\cdots\]
	is $\mathbb S_{(p)}$.
\end{theorem}
Furthermore, there are pullback squares
% https://q.uiver.app/#q=WzAsNCxbMCwwLCJMX25cXG1hdGhiYiBTX3socCl9Il0sWzAsMSwiTF97bi0xfVxcbWF0aGJiIFNfeyhwKX0iXSxbMSwwLCJMX3tLKG4pfVxcbWF0aGJiIFNfeyhwKX0iXSxbMSwxLCJMX3tLKG4tMSl9XFxtYXRoYmIgU197KHApfSJdLFsxLDNdLFswLDFdLFswLDJdLFsyLDNdXQ==&macro_url=https%3A%2F%2Fraw.githubusercontent.com%2FdFoiler%2Fnotes%2Fmaster%2Fnir.tex
\[\begin{tikzcd}[cramped]
	{L_n\mathbb S_{(p)}} & {L_{K(n)}\mathbb S_{(p)}} \\
	{L_{n-1}\mathbb S_{(p)}} & {L_{K(n-1)}\mathbb S_{(p)}}
	\arrow[from=1-1, to=1-2]
	\arrow[from=1-1, to=2-1]
	\arrow[from=1-2, to=2-2]
	\arrow[from=2-1, to=2-2]
\end{tikzcd}\]
which allows us to understand $\mathbb S_{(p)}$ from its completions $L_{K(a_1)}\cdots L_{K(a_n)}\mathbb S$.
\begin{remark}
	This pullback square is analogous to the following one.
	% https://q.uiver.app/#q=WzAsNCxbMCwwLCJcXFpaX3socCl9Il0sWzEsMCwiXFxaWl97KHApfV5cXGxhbmQiXSxbMSwxLCJcXFFRX3AiXSxbMCwxLCJcXFFRIl0sWzAsMywiIiwwLHsic3R5bGUiOnsidGFpbCI6eyJuYW1lIjoiaG9vayIsInNpZGUiOiJ0b3AifX19XSxbMywyLCIiLDAseyJzdHlsZSI6eyJ0YWlsIjp7Im5hbWUiOiJob29rIiwic2lkZSI6InRvcCJ9fX1dLFswLDEsIiIsMix7InN0eWxlIjp7InRhaWwiOnsibmFtZSI6Imhvb2siLCJzaWRlIjoidG9wIn19fV0sWzEsMiwiIiwyLHsic3R5bGUiOnsidGFpbCI6eyJuYW1lIjoiaG9vayIsInNpZGUiOiJ0b3AifX19XSxbMCwyLCIiLDEseyJzdHlsZSI6eyJuYW1lIjoiY29ybmVyIn19XV0=&macro_url=https%3A%2F%2Fraw.githubusercontent.com%2FdFoiler%2Fnotes%2Fmaster%2Fnir.tex
	\[\begin{tikzcd}[cramped]
		{\ZZ_{(p)}} & {\ZZ_{(p)}^\land} \\
		\QQ & {\QQ_p}
		\arrow[hook, from=1-1, to=1-2]
		\arrow[hook, from=1-1, to=2-1]
		\arrow["\lrcorner"{anchor=center, pos=0.125}, draw=none, from=1-1, to=2-2]
		\arrow[hook, from=1-2, to=2-2]
		\arrow[hook, from=2-1, to=2-2]
	\end{tikzcd}\]
\end{remark}
Thus, for a given height $h$, we would like to understand the localizations and how they glue together. For a single completion, one has Galois theory.
\begin{theorem}
	For each height $h$, set $k\coloneqq\FF_{p^h}$, and let $\Gamma$ be the Honda formal group of height $h$, and set $G_h\coloneqq\op{Aut}\Gamma$. Then there is an equivalence
	\[L_{K(h)}\mathbb S\to E(k,\Gamma)^{hG_h\times\op{Gal}(k/\FF_p)}.\]
	In other words, the category of $K(h)$-local $\mathbb S$-modules is equivalent to the category of $G_h\rtimes\op{Gal}(k/\FF_p)$-equivariant $(h)$-local $E(k,\Gamma)$-modules.
\end{theorem}
\begin{remark}
	One can use this to prove that $L_n$ is smashing.
\end{remark}
\begin{remark}
	The homotopy groups of $L_{K(n)}\mathbb S$ can be computed from the spectral sequence
	\[\mathrm H^*(G_h\rtimes\op{Gal}(k/\FF_p);W(k)[[u_1,\ldots,u_{h-1}]][u,1/u])\Rightarrow\pi_*L_{K(n)}\mathbb S.\]
	When $h=1$, this is related to an $L$-function. We remark that the domain can be extended from $\ZZ$ to $K(h)$-local spectra $V$ admitting some $W$ for which $L_{K(h)}(V\otimes W)=L_{K(h)}\mathbb S$; the functional equation for the relevant $L$-function actually knows about ``Gross--Hopkins duality'' of the noninteger elements of the domain.
\end{remark}
Now that we have some understanding of one completion, one can try to understand multiple completions. This is handled by the following conjecture.
\begin{conj}[Hopkins] \label{conj:hopkins}
	For $h<p$,
	\[L_{h-1}L_{K(h)}\mathbb S_{(p)}\stackrel?\cong\bigoplus_{\substack{0\le j\le j\\1<i_1<\cdots<i_j\le h}}\bigotimes_{k=1}^jL_{h-i_k}\mathbb S^{1-2i_k}.\]
\end{conj}
\begin{remark}
	This is known for $h\le2$, except at $(p,h)=(2,2)$
\end{remark}
However, this conjecture is false in general for $h<p$.
\begin{remark}
	This conjecture would imply that
	\[L_{K(h-1)}L_{K(h)}\mathbb S=L_{K(h-1)}\mathbb S\oplus\Sigma^{-1}L_{K(h-1)}\mathbb S.\]
	Indeed, it was recently shown that
	\[L_{K(1)}L_{K(2)}L_{K(3)}\mathbb S_{(2)}=L_{K(1)}\mathbb S_{(2)}\oplus\Sigma^{-1}L_{K(1)}\mathbb S_{(2)}\oplus\Sigma^{-2}L_{K(1)}\mathbb S/2\oplus\Sigma^{-3}L_{K(1)}\mathbb S/2.\]
	We remark that the right-hand side of \Cref{conj:hopkins} is strictly smaller than the above right-hand side.
\end{remark}
\begin{remark}
	The conjecture would imply that $\mathbb S^\land_p$ is a direct summand of $\prod_{h\ge1}L_{K(h)}\mathbb S$, which remains open.
\end{remark}

\subsection{The Chromatic Splitting Conjecture}
The target of the course was the following corollary of \Cref{conj:hopkins}.
\begin{theorem}[Barthel--Mann--Schlank--Senger--Weinstein--Zhou] \label{thm:target}
	Fix a height $h$. If $(p-1)\nmid(h-1)$, then $L_{K(h-1)}L_{K(h)}\mathbb S$ has $L_{K(h-1)}\mathbb S\oplus\Sigma^{-1}L_{K(h-1)}\mathbb S$ as a direct summand.
\end{theorem}
\begin{remark}
	One expects to have equality, not merely a direct summand.
\end{remark}
\begin{remark}
	This was proven without computing $L_{K(h-1)}L_{K(h)}\mathbb S$.
\end{remark}
\begin{proof}[Sketch]
	Let $L$ be the Torii extension of $k$, and we let $\widehat L$ be its completion, which is perfectoid. Then there is a composite
	\[E(k,\Gamma_{h-1})\to E(L,\Gamma_h)\to E(\widehat L,\Gamma_h).\]
	The moral is that $E(L,\Gamma_h)$ should be related to $L_{K(h)}\mathbb S$, so it is impressive that we have managed to produce a map out of this object. Now, we showed last time that
	\[E(k,\Gamma_{h-1})\cong\left(E(\widehat L,\Gamma_h)^{hG_h^\circ}\right)^{h\FF_p^\times},\]
	where $G_h^\circ$ is the kernel of $\det\colon G_h\to\ZZ_p^\times$. (This equivalence requires $(p-1)\nmid(h-1)$.) In fact, this is an equivalence of $G_{h-1}\rtimes\op{Gal}(k/\FF_p)$-equivariant objects, and taking fixed points produces the desired result.
\end{proof}\begin{remark}
	One can try to retell the whole story for $L$ instead of $\widehat L$. Instead of having a perfect ring, we now have a Gelfand ring, and Scholze's school is trying to figure out what the relevant Dieudonn\'e theory looks like.
\end{remark}
There is still some hope of proving a theorem when $(p-1)\mid(h-1)$. Without any hypothesis on $p$ and $h$, we successfully computed
\[\mathrm H^i_{\mathrm{cts}}\left(G_h^\circ;\widehat L\right)=\mathrm H^i\left(\op{BC}\left(-\frac{1}{h-1}\right)^*;\OO^\flat\right)=\begin{cases}
	k & \text{if }i\in\{0,h\}, \\
	0 & \text{otherwise}.
\end{cases}\]
There is a residual $\FF_p^\times$-action, a residual $(1+p\ZZ_p)$-action, and a residual $G_{h-1}\rtimes\op{Gal}(k/\FF_p)$-action. To prove \Cref{thm:target}, we only used the action by $\FF_p^\times$ to kill the cohomology in degree $i=h$. To do better, one can try to use the other actions; for example, $G_{h-1}\rtimes\op{Gal}(k/\FF_p)$ is given by taking to determinant to $\ZZ_p^\times$, which then surjects onto $\FF_p^\times$ and embeds to $k^\times$.

Thus, we may consider the cofiber
\begin{equation}
	E(k,\Gamma_{h-1})\to\left(E(\widehat L,\Gamma_h)^{hG_h^\circ}\right)^{h\FF_p^\times}\to C. \label{eq:chromatic-splitting-small-primes}
\end{equation}
We know that $C$ is a $K(h-1)$-local $G_{h-1}$-equivariant $E(k,\Gamma_{h-1})$-module. Furthermore, its homotopy group is at degree $h$, where it is some $k$ with some action, so it follows that $C^{hG_{h-1}\rtimes\op{Gal}(k/\FF_p)}$ is a $K(h-1)$-local Picard element. Here are some remaining questions.
\begin{ques}
	What is the $K(h-1)$-local Picard element $C^{hG_{h-1}\rtimes\op{Gal}(k/\FF_p)}$?
\end{ques}
\begin{ques}
	Does the cofiber sequence \eqref{eq:chromatic-splitting-small-primes} split?
\end{ques}
Here is a purely algebraic problem, which is a place to start.
\begin{ques}
	The module
	\[\mathrm R\Gamma\left(\op{BC}\left(-\frac1{h-1}\right)^*;\OO^\flat\right)\]
	is a condensed $\ZZ$-module with $G_{h-1}$-action. Is it a direct sum of $k$ and $\Sigma^hk(1)$?
\end{ques}
We close our course with another conjecture.
\begin{conj}[Deperfection] \label{conj:deperfection}
	The natural map
	\[E(L,\Gamma)^{h(G_h\times G_{h-1})\rtimes\op{Gal}(k/\FF_p)}\to E(\widehat L,\Gamma)^{h(G_h\times G_{h-1})\rtimes\op{Gal}(k/\FF_p)}\]
	is an equivalence.
\end{conj}
\begin{remark}
	This is known for $h=2$, even when $p=2$.
\end{remark}
\begin{remark}
	\Cref{conj:deperfection} admits an algebraic reformulation. Recall that $K=k((u_{h-1}))$, and $L$ is some Galois extension of $K$. Now, $\widehat K=k((u^{1/p^\infty}))$, and this isomorphism is actually meaningful now. \Cref{conj:deperfection} is then equivalent to
	\[\mathrm H^*_{\mathrm{cts}}\left(G_h;\widehat K/K\right)=0.\]
	Indeed, this implies that the difference between $\widehat L$ and $L$ is zero. In short, $E(\widehat L)^{hG_{h-1}}=E(\widehat K)$.
\end{remark}
\begin{example}
	When $h=2$, one finds that
	\[G_2=\left(\frac{W(\FF_{p^2})\langle S\rangle}{\left(S^2-p,a^\sigma S-Sa\right)}\right)^\times,\]
	and $G_2$ acts on $\pi_0(E(k,\Gamma_2)/p)=W(\FF_{p^2})/p[[u_1]]$. One can use a computer to compute the relevant action and check that \Cref{conj:deperfection} is true for $(p,h)=(5,2)$.
\end{example}

\end{document}